\documentclass[final]{beamer}  
  \mode<presentation> { \usetheme{ParisSaclay}  }
  \usepackage[spanish]{babel}
  \usepackage[utf8]{inputenc}
  \usepackage{amsmath,amsthm, amssymb, latexsym}

  \boldmath %math en gras
  \usepackage[orientation=portrait,size=a1,scale=1.4]{beamerposter}

  \usepackage{graphicx}
  \usepackage{tikz}
  
  \title{Integración de sentimiento de mercado en agentes de RL para la optimización de trading}
  \author{Julian Duarte \quad Julian Jimenez}
  \institute[]{Ciencia de Datos -- Departamento de Matemáticas \\ \medskip Tutores: Prof. Julián Fernando Sánchez \quad Prof. Juan Carlos Urueña \\ \medskip \textbf{Semillero de Aprendizaje por Refuerzo}}

\begin{document}

\begin{frame} [fragile]{} 

\begin{columns}
    \begin{column}{.46\textwidth}
        
    \begin{block}{1. Introducción: ¿Por qué una IA que sienta el mercado?}
        El mercado cripto opera como una ``montaña rusa'' impulsada por la psicología colectiva. Sesgos cognitivos como el \textbf{pánico sistémico} o la \textbf{euforia irracional} (FOMO) distorsionan la toma de decisiones, creando ineficiencias críticas.
        \begin{itemize}
            \item \textbf{Problema:} La falta de un "termómetro emocional" en los algoritmos convencionales.
            \item \textbf{Propuesta:} Un agente racional entrenado mediante \textit{RL} para navegar la volatilidad integrando el sentimiento humano.
        \end{itemize}
        \centering
        \vspace{0.3cm}
        \includegraphics[width=0.85\linewidth]{ciclo_psicologico_mercado.png}
    \end{block}

    \vspace{0.5cm}

    \begin{block}{2. Objetivos del Proyecto}
        \begin{itemize}
            \item \textbf{Política Óptima ($\pi^*$):} Converger a un manual de reglas de decisión que maximice el beneficio acumulado ($G_t$).
            \item \textbf{Fusión de Datos:} Integrar el índice \textit{Fear \& Greed} como una variable de estado exógena crítica.
            \item \textbf{Rentabilidad Eficiente:} Priorizar la búsqueda de retornos con baja volatilidad (Sharpe Ratio optimizado).
        \end{itemize}
    \end{block}

    \vspace{0.5cm}

    \begin{alertblock}{3. Metodología: El Tablero de Instrumentos}
        Para decidir, el agente procesa un vector de estado ($S_t$) con cuatro ``sensores'' fundamentales que resumen el entorno:
        \begin{itemize}
            \item \textbf{Dirección de Mercado ($M_t$):} Indica el régimen de precio actual (¿el mercado tiene tendencia alcista o bajista?).
            \item \textbf{Termómetro Emocional ($Sent_t$):} Valor del índice \textit{Fear \& Greed}, que alerta al agente sobre pánicos o euforias colectivas.
            \item \textbf{Sensor de Fatiga ($RSI_t$):} Indica si el precio está ``agotado'' (sobrecompra o sobreventa), ayudando a evitar entradas tardías.
            \item \textbf{Memoria de Posición ($I_t$):} Informa al agente si posee activos (BNB) o efectivo (USDT), condicionando sus acciones posibles.
        \end{itemize}
        \textbf{Aprendizaje (Q-Learning):}
        Basándose en estos datos, el agente elige una acción $\{ \text{Compra, Venta, Espera} \}$ y optimiza su política mediante el error de predicción temporal:
        \[ Q(s, a) \leftarrow Q(s, a) + \alpha [r + \gamma \max_{a'} Q(s', a') - Q(s, a)] \]
        \centering
        \vspace{0.3cm}
        \includegraphics[width=0.85\linewidth]{grafica_montana_rusa.png}
    \end{alertblock}
    
    \end{column}

    \begin{column}{.46\textwidth}	

    \begin{block}{4. Evaluación: Midiendo la sabiduría de la IA}
        Evaluamos la resiliencia algorítmica mediante métricas estándar de riesgo-retorno:
        \begin{itemize}
            \item \textbf{Sharpe Ratio:} Medición de la rentabilidad eficiente (ganar de forma tranquila y constante).
            \item \textbf{Máximo Drawdown:} El termómetro del "peor momento" para proteger el capital ante caídas abruptas.
            \item \textbf{Generalización Out-of-Sample:} Validación de la política $\pi$ en periodos de tiempo no vistos durante el entrenamiento.
        \end{itemize}
        \centering
        \vspace{0.3cm}
        \includegraphics[width=0.85\linewidth]{grafica_comparativa_riesgo.png}
    \end{block}
    
    \vspace{0.5cm}	

    \begin{exampleblock}{5. Conclusiones y Aplicación}
        La investigación demuestra que una arquitectura resiliente es aquella capaz de interpretar el pulso emocional del mercado. El modelo propuesto actúa como un \textbf{freno racional} teórico ante la impulsividad humana.

        \textbf{Aportes y Proyecciones:}
        \begin{itemize}
            \item \textbf{Resiliencia Algorítmica:} La integración de finanzas conductuales permite al agente modelar comportamientos defensivos ante pánicos colectivos.
            \item \textbf{Soporte de Decisión:} El marco conceptual sirve como base para sistemas que validen la racionalidad de una inversión antes de su ejecución.
            \item \textbf{Fase de Implementación:} Se proyecta el uso de \textit{Websockets} para la transición de simulaciones hacia entornos de mercado real en Binance.
            \item \textbf{Escalabilidad:} La metodología es aplicable a la gestión de portafolios multi-activo y al análisis de sentimiento mediante NLP.
        \end{itemize}
    \end{exampleblock}

    \vspace{0.5cm}

    \begin{block}{6. Contacto y Referencias}
        \small
        Sutton \& Barto (2018). \textit{Reinforcement Learning: An Introduction}. MIT Press.
        \smallskip
        \centering
        \faEnvelope \quad julian.duarte2 / julian.jimenez5 @est.uexternado.edu.co \\
        \textbf{Semillero de Aprendizaje por Refuerzo - 2026}
        \vfill
        \centering
        \small \textit{Escanea el código QR para acceder al repositorio y documentación completa del proyecto.}
    \end{block}

    \end{column}
\end{columns}

\end{frame}

\end{document}